\documentclass[10pt, journal]{IEEEtran}

% --- Modular Packages & Preamble Manifest ---
\input{packages.tex}

% --- Custom Colors & Callout Blocks ---
\definecolor{ieeeblue}{HTML}{006699}
\definecolor{lumabg}{HTML}{F4F4F4}
\definecolor{blockquoteleft}{HTML}{777777}
\definecolor{blockquotebg}{HTML}{EAEAEA}

% Code Listing Preferences
\lstset{
  backgroundcolor=\color{lumabg},
  basicstyle=\ttfamily\scriptsize, 
  breaklines=true,
  breakatwhitespace=true,
  frame=leftline,
  framerule=2pt,
  rulecolor=\color{gray},
  framesep=6pt,
  xleftmargin=6pt,
  showstringspaces=false,
  language=Python
}

% --- Document Metadata ---
\title{__TITLE__}
\author{__AUTHOR__,~\IEEEmembership{Member,~IEEE}%
\thanks{The author is with the __DEPENDENCY__, Delft University of Technology, 2628 CD Delft, The Netherlands (e-mail: author@tudelft.nl).}}

\begin{document}

\maketitle

% --- Abstract & Keywords Queue ---
\begin{abstract}
Autonomous robotic swarms offer transformative potential in subterranean exploration, disaster response, and planetary speleology. In these scenarios, the absence of global navigation satellite systems (GNSS) presents a fundamental challenge. Individual agents equipped solely with visual-inertial odometry (VIO) inevitably experience unbounded dead-reckoning drift over prolonged mission horizons. We formulate a decentralized factor graph framework with inter-agent ultra-wideband ranging constraints to bound global localization drift.
\end{abstract}

\begin{IEEEkeywords}
Decentralized Estimation, Factor Graphs, Autonomous Swarms, Ultra-Wideband, Robotic Networks.
\end{IEEEkeywords}

% --- Section 1: Introduction ---
\section{Introduction and Problem Formulation}
\IEEEPARstart{T}{he} primary objective of an IEEE journal paper is to convey technical engineering concepts within a highly uniform, space-efficient, two-column layout grid. To maintain this uniformity, authors must strictly adhere to a standardized hierarchy of structural boundaries.

\subsection{Kinematic State Propagation}
Consider a swarm comprising $N$ autonomous agents. Each agent $i$ tracks its state vector $\mathbf{x}_i$:
\begin{equation}
	\mathbf{x}_i = \begin{bmatrix} \mathbf{p}_i \\ \mathbf{R}_i \\ \mathbf{v}_i \end{bmatrix}
\end{equation}

\section{Distributed Consensus Optimization}
Equations inside IEEE papers run inside centered blocks and are numbered sequentially along the right-hand margin:
\begin{align}
	\dot{\mathbf{x}}(t) &= \mathbf{A}\mathbf{x}(t) + \mathbf{B}\mathbf{u}(t) \label{eq:statespace} \\
	\mathbf{y}(t)       &= \mathbf{C}\mathbf{x}(t) + \mathbf{D}\mathbf{u}(t) \label{eq:output}
\end{align}

\begin{lstlisting}[language=Python, caption={Consensus gradient descent loop.}]
def compute_step(x, neighbors, weights):
    grad = sum(w * (x - x_n) for x_n, w in zip(neighbors, weights))
    return x - 0.05 * grad
\end{lstlisting}

\section{Citations \& Scholarly References}
Citations map directly to the bibliography file \cite{ieee_evans, ieee_ramirez}.

\bibliographystyle{IEEEtran}
\bibliography{library}
\end{document}
